% Solution: Tadelis (2013) Exercise 9.1 "Multistage Practice". Input by \solitem.
% Reused from lista-reais-lectures-01-04-solutions (item 19), deepened. Check: check_C_L05.py.
\begin{exsol}{Tadelis Exercise 9.1}{186--187}{Exam level}{4}
\asked For the $3\times3$ game played twice with the first-period outcome observed: (a) all
pure-strategy SPE with $\delta = 1$, with history-contingent strategies; (b) for each, the smallest
$\delta$ that supports it.

\method Multistage games (Tadelis \S9.1--9.4, pp.~175--184) and the finitely repeated game
$G(2,\delta)$ (\S10.1, pp.~190--192). Period-2 play must be a stage Nash equilibrium after every
history (Prop.~9.2, p.~181). A first-period profile that is not a stage equilibrium can be
supported only by a carrot (a good stage equilibrium on the path) and a stick (a bad one after a
deviation), and only when $\delta \ge \delta^* = g/(c - s)$ (\texttt{Notas/L04/L04-02},
\texttt{Notas/L05/L05-02}; Tadelis's own two-stage example, p.~191). Deviations are checked one
stage at a time (one-stage deviation principle, Thm.~9.1, p.~185). Slides: Lecture~4, slide~22
($G(T,\delta)$, payoff $\sum_t \delta^{t-1}u_i^t$, each period starts a subgame) and slide~17 (SPE).

\solution
\textbf{Notation.} $v_i = u_i(a^1) + \delta\, u_i(a^2)$, where $a^t$ is the period-$t$ action
profile. A pure strategy for player~$i$ is a period-1 action and a map from the nine possible
period-1 outcomes to a period-2 action. $g_i(a)$ is player~$i$'s best one-shot gain from deviating
from $a$:
\[
  g_1(a) = \max_{r} u_1(r, a_2) - u_1(a), \qquad g_2(a) = \max_{c} u_2(a_1, c) - u_2(a).
\]
Best replies are underlined:
\begin{center}
\begin{tabular}{r|c|c|c|}
\multicolumn{1}{r}{} & \multicolumn{1}{c}{$L$} & \multicolumn{1}{c}{$C$} & \multicolumn{1}{c}{$R$}\\ \cline{2-4}
$T$ & $10,\ 10$ & $2,\ 12$ & $0,\ \br{13}$ \\ \cline{2-4}
$M$ & $12,\ 2$ & $\br{5},\ \br{5}$ & $0,\ 0$ \\ \cline{2-4}
$B$ & $\br{13},\ 0$ & $0,\ 0$ & $\br{1},\ \br{1}$ \\ \cline{2-4}
\end{tabular}
\end{center}

\textbf{Step 1: stage equilibria.} Player~1's best replies: to $L$, $B$ (13); to $C$, $M$ (5); to
$R$, $B$ (1). Player~2's: to $T$, $R$ (13); to $M$, $C$ (5); to $B$, $R$ (1). The mutual best
replies are
\[
  (M, C)\ \text{with } (5,5) \quad\text{(the carrot)}, \qquad (B, R)\ \text{with } (1,1)
  \quad\text{(the stick)} .
\]
Each player gets 5 in the carrot and 1 in the stick, so the largest punishment one period can
deliver is $c - s = 5 - 1 = 4$.

\textbf{Step 2: period 2 (the last phase).} After any first-period outcome the subgame is the stage
game (plus a sunk constant), so it must be $(M,C)$ or $(B,R)$. No one-stage deviation in period~2
pays, because both are stage equilibria and nothing follows.

\textbf{Step 3: period 1 (which profiles can be supported).} A first-period profile $a$ is
supportable if and only if each player's gain is at most what the threat can take away:
\begin{equation}
  g_i(a) \;\le\; \delta\,(5 - 1) = 4\delta, \qquad i = 1, 2, \label{eq:tadelis-9-1-ic}
\end{equation}
using $(M,C)$ after $a$ and $(B,R)$ after any deviation. No continuation can do better, since 5 is
the most and 1 the least a player can get in period~2. For a stage equilibrium $a$, $g_i(a) = 0$
and any continuation works. The gains are:
\begin{center}
\small
\begin{tabular}{lccccc}
\toprule
$a$ & $u(a)$ & $g_1(a)$ & $g_2(a)$ & supportable ($\delta=1$)? & $\delta^* = \max_i g_i/4$\\
\midrule
$(T,L)$ & $(10,10)$ & $13-10=3$ & $13-10=3$ & yes & $3/4$ \\
$(T,C)$ & $(2,12)$ & $5-2=3$ & $13-12=1$ & yes & $3/4$ \\
$(T,R)$ & $(0,13)$ & $1-0=1$ & $0$ & yes & $1/4$ \\
$(M,L)$ & $(12,2)$ & $13-12=1$ & $5-2=3$ & yes & $3/4$ \\
$(M,C)$ & $(5,5)$ & $0$ & $0$ & yes & $0$ \\
$(M,R)$ & $(0,0)$ & $1-0=1$ & $5-0=5$ & no & -- \\
$(B,L)$ & $(13,0)$ & $0$ & $1-0=1$ & yes & $1/4$ \\
$(B,C)$ & $(0,0)$ & $5-0=5$ & $1-0=1$ & no & -- \\
$(B,R)$ & $(1,1)$ & $0$ & $0$ & yes & $0$ \\
\bottomrule
\end{tabular}
\end{center}

\textbf{(a) The pure SPE with $\delta = 1$.} There are nine SPE \emph{paths}:
\begin{itemize}
  \item the four sequences of stage equilibria, played regardless of history (Prop.~9.1, p.~180):
  $(M,C)$ then $(M,C)$; $(M,C)$ then $(B,R)$; $(B,R)$ then $(M,C)$; $(B,R)$ then $(B,R)$;
  \item the five carrot-and-stick paths: $a$ then $(M,C)$, for $a \in \{(T,L), (T,C), (T,R),
  (M,L), (B,L)\}$.
\end{itemize}
For a carrot-and-stick path the strategies are, for $\hat a = (\hat a_1, \hat a_2)$ one of those five:
\begin{align*}
  s_1&:\ \text{play } \hat a_1;\ \text{in period 2 play } M \text{ if } \hat a \text{ was played, } B
  \text{ otherwise};\\
  s_2&:\ \text{play } \hat a_2;\ \text{in period 2 play } C \text{ if } \hat a \text{ was played, } R
  \text{ otherwise}.
\end{align*}
\emph{One-stage deviation checks for each path} (with $\delta = 1$; path payoff against best deviation):
\begin{center}
\small
\begin{tabular}{lcccc}
\toprule
$\hat a$ & P1: path & P1: best deviation & P2: path & P2: best deviation\\
\midrule
$(T,L)$ & $10+5=15$ & $B$: $13+1=14$ & $10+5=15$ & $R$: $13+1=14$\\
$(T,C)$ & $2+5=7$ & $M$: $5+1=6$ & $12+5=17$ & $R$: $13+1=14$\\
$(T,R)$ & $0+5=5$ & $B$: $1+1=2$ & $13+5=18$ & none (best reply)\\
$(M,L)$ & $12+5=17$ & $B$: $13+1=14$ & $2+5=7$ & $C$: $5+1=6$\\
$(B,L)$ & $13+5=18$ & none (best reply) & $0+5=5$ & $R$: $1+1=2$\\
\bottomrule
\end{tabular}
\end{center}
In period~2 both players play a stage equilibrium after every history, so no deviation pays there.
$(M,R)$ and $(B,C)$ cannot be supported: a gain of 5 exceeds the largest punishment, 4. For instance,
at $(M,R)$ player~2 gets $0 + 5 = 5$ on the path but $5 + 1 = 6$ by deviating to $C$.

These strategies are not the only ones. Any assignment of $(M,C)$ or $(B,R)$ to the nine period-1
outcomes that satisfies \eqref{eq:tadelis-9-1-ic} for the deviations that would be profitable also
works (e.g.\ for $(T,R)$ only player~1's deviation to $B$ must be punished). Counting every such map,
\texttt{check\_C\_L05.py} finds 992 pure SPE strategy profiles at $\delta = 1$, all with one of the nine
paths above.

\textbf{(b) The smallest discount factor.} From \eqref{eq:tadelis-9-1-ic},
$\delta^*(a) = \max_i g_i(a)/4$ (last column of the first table). Solving each binding inequality:
\begin{itemize}
  \item $(T,L)$: $10 + 5\delta \ge 13 + \delta \iff 4\delta \ge 3 \iff \delta \ge \frac34$ (both players);
  \item $(T,C)$: player~1 binds, $2 + 5\delta \ge 5 + \delta \iff \delta \ge \frac34$; player~2 needs
  $12 + 5\delta \ge 13 + \delta \iff \delta \ge \frac14$;
  \item $(M,L)$: the mirror image of $(T,C)$, $\delta \ge \frac34$;
  \item $(T,R)$: only player~1, $0 + 5\delta \ge 1 + \delta \iff \delta \ge \frac14$;
  \item $(B,L)$: only player~2, $0 + 5\delta \ge 1 + \delta \iff \delta \ge \frac14$;
  \item the four sequences of stage equilibria: any $\delta \ge 0$.
\end{itemize}
Below $\delta^*$ no continuation map works, because $(M,C)$ after $\hat a$ and $(B,R)$ after the
deviation is already the largest reward gap available.

\textbf{A caveat: renegotiation.} The stick $(B,R)$ is Pareto-dominated by $(M,C)$. After a
deviation both players would like to ``forgive'' and play $(M,C)$ in period~2, and if they can
renegotiate the threat fails (Gibbons \S2.3.A, pp.~87--88; \texttt{Notas/L05/L05-02}). The SPE
above take the threat as given.

\interp Two stage equilibria with different payoffs give the second period a carrot and a stick.
That lets the players reach $(T,L)$, which pays $(10,10)$, twice the $(5,5)$ of the good stage
equilibrium, in the first period. Profiles with a gain above $4$ cannot be enforced by a single
future period, however patient the players are. This is the finite-horizon version of the logic
of Lecture~5: what can be sustained today is limited by how much the future can reward and punish.
With a unique stage equilibrium (item Cabral~7.13) that gap is zero and nothing can be sustained.

\verify \texttt{check\_C\_L05.py} finds the two stage equilibria and enumerates all
$9 \times 2^9$ combinations of first-period profile and continuation map at $\delta = 1$ (992 SPE,
9 paths; $(M,R)$ and $(B,C)$ excluded). For every threshold it checks that each path is an SPE at
$\delta^*$ and at $\delta^* + \frac1{1000}$, that no continuation map at all supports it at $\delta^* -
\frac1{1000}$, and that on the grid $\delta = k/1000$ the carrot-and-stick map works exactly when
$\delta \ge \delta^*$. Output: ``Tadelis 9.1: stage NE (M,C),(B,R); 992 pure SPE at delta=1 on 9 paths;
delta* = 3/4 (TL, TC, ML), 1/4 (TR, BL); grid and +-1/1000 pass''. Tadelis prints no answer.

\mistakes
\begin{itemize}
  \item Listing only the outcome ``$(T,L)$ then $(M,C)$'' without the off-path play. The question asks
  for history-contingent strategies, and the threat to play $(B,R)$ after a deviation is the whole
  point.
  \item Using the smallest gain instead of the largest: the threshold is $\max_i g_i/4$, so in $(T,C)$
  player~1's gain of 3 binds, not player~2's gain of 1.
  \item Claiming $(M,R)$ or $(B,C)$ can be sustained with enough patience. With one period left the
  punishment is capped at $5 - 1 = 4$, which is below the gain of $5$ even at $\delta = 1$.
\end{itemize}

\answer{Stage equilibria: $(M,C)$ $(5,5)$ and $(B,R)$ $(1,1)$. (a) Pure SPE paths with $\delta=1$:
the four sequences of stage equilibria, and $a$ then $(M,C)$ for $a \in \{(T,L),(T,C),(T,R),
(M,L),(B,L)\}$, supported by ``play $(M,C)$ after $a$, $(B,R)$ after any other outcome''.
$(M,R)$ and $(B,C)$ cannot be supported. (b) $\delta^* = \frac34$ for $(T,L)$, $(T,C)$, $(M,L)$;
$\frac14$ for $(T,R)$, $(B,L)$; $0$ for the sequences of stage equilibria.}
\end{exsol}
